\section{SAT, SMT, and first-order proving}
\label{sec:automated-reasoning-06}

\begin{definition}[Satisfiability modulo theories]
Given a specified background theory \(T\), the \emph{satisfiability
modulo theories problem} asks whether some model of \(T\) satisfies the
input formula under some variable assignment; equivalently, whether
\(T\) together with the existential closure of the input formula is
satisfiable.
\end{definition}

An SMT solver commonly combines Boolean search with decision procedures for
fragments such as equality with uninterpreted functions, linear arithmetic,
arrays, or bit-vectors.  The exact input grammar, sort system, theory symbols,
commands, and responses are standardized by SMT-LIB.  Version 2.7 specifies a
many-sorted language, operational commands, and logical semantics
\cite{SMTLIB2025}.  A response of \texttt{sat} should be accompanied by a model
when one is needed; a response of \texttt{unsat} should be accompanied by a
checkable proof or trusted-solver assumption when the conclusion is used in a
high-assurance development.

The TPTP language supplies standardized formula roles, names, syntax, and
annotations for automated theorem proving \cite{TPTP}.  A parser-accepted file
is not thereby a correctly formalized theory.  The declarations, formula
roles, conjecture polarity, and imported axioms must still be audited.

